\subsection{因子整环上的多项式环}

设 A 为因子整环，K 为其分式域，f 为 K{[}X{]} 中的非零元素；K* 中的元素 c 若是 f 的系数的最大公因元，则称为 f 的容度。设 v 是 K 上对 A 为本性的赋值，\(\bar{v}\) 是它到 K{[}X{]} 的典范延拓（定义为 \(\bar{v}\left(\sum_{i}a_{i}X^{i}\right)=\inf v(a_{i})\)；参见第VI章第10节第1号命题2）；则 \(\bar{v}(f)=v(c)\)。

引理1（高斯）. 设 \(f, f'\) 是 \(\mathbf{K}[\mathbf{X}]\) 中的非零元素，\(c, c'\) 是 \(f, f'\) 的容度。则 \(cc'\) 是 \(ff'\) 的一个容度。

设 \(d\) 是 \(ff'\) 的容度。对于每个赋值 \(v\)，它定义在 \(K\) 上且对 \(A\) 为本性，令 \(\bar{v}\) 表示它到 \(K[X]\) 的典范延拓。则

\[
v (d) = \bar {v} (f f ^ {\prime}) = \bar {v} (f) + \bar {v} (f ^ {\prime}) = v (c) + v (c ^ {\prime}) = v (c c ^ {\prime}).
\]

因此 \(cc'd^{-1}\) 是 A 的可逆元素。

定理2. 设 A 为因子整环，K 为其分式域，\((p_{\iota})\) 为 A 的极值元代表系，\((\mathrm{P}_{\lambda})\) 为 K{[}X{]} 中不可约多项式的代表系，且每个 \(P_{\lambda}\) 的容度为1。则：

\begin{enumerate}
\def\labelenumi{(\roman{enumi})}
\item
  A{[}X{]} 是因子整环；
\item
  \(p_{\iota}\) 和 \(P_{\lambda}\) 的集合是 A{[}X{]} 的极值元代表系。
\end{enumerate}

设 \(f\) 是 \(A[X]\) 中的非零元素。在环 \(K[X]f\) 中，f 可以唯一分解为如下形式：

\[
f = a \prod_ {\lambda} \mathrm{P} _ {\lambda} ^ {n (\lambda)} \quad (a \in \mathrm{K} ^ {*}, n (\lambda) \geqslant 0).
\]

引理1表明 \(a\) 是 \(f\) 的容度。因此 \(a \in \mathbf{A}\)。由于 \(\mathbf{A}\) 是因子整环，\(a\) 可以唯一分解为如下形式：

\[
a = u \prod_ {\iota} p _ {\iota} ^ {m (\iota)} \quad (u \text {   invertible   in   A }, m (\iota) \geqslant 0).
\]

于是得到以下分解的存在性和唯一性：

\[
f = u \prod_ {\iota} p _ {\iota} ^ {m (\iota)} \prod_ {\lambda} \mathrm{P} _ {\lambda} ^ {n (\lambda)}.
\]

注意，该定理证明了 A 的每个元素在 A 和 A{[}X{]} 中都有相同的极值元分解。因此，A 的一族元素的最大公因元在 A 和 A{[}X{]} 中相同。

也可以利用第1节第10号命题18证明，A{[}X{]} 是因子整环当且仅当 A 是因子整环。

推论。若 A 是因子整环，则环 A{[}X \(_{1}\) , . . . , X \(_{n}\) {]} 是因子整环。

按 n 归纳论证。

该推论可推广到具有无限族不定元的情形（参见习题2）。

